\chapter{Brief History}

On September the 8th 1930, David Hilbert gives his final address before retirement in K\"onigsberg, Germany. Passionately, he insists future scientists to resist the notion of the \emph{Ignorabimus}; the idea that certain truths will forever out of humanity's reach:

\begin{displayquote}
We must not believe those, who today with philosophical bearing and a tone of superiority prophesy the downfall of culture and accept the ignorabimus. 

For us there is no ignorabimus, and in my opinion even none whatever in natural science. In place of the foolish of ignorabimus let stand our slogan:

{\centering\emph{We must know. We will know\dots}\par}
{\raggedleft --- David Hilbert (1930, tr. James T. Smith)\par}
\end{displayquote}

A day earlier Kurt G\"odel had quietly announced his Incompleteness Theorems, proving that certain truths will be forever out of humanity's reach. To some, this moment marks the end of the decades work that was Hilbert's Program, a project he had initiated to save mathematics from the foundational crisis that it was in.

%  Bernhard Riemann 1826–1866
% Turning Points in the Conception of Mathematics (Contains distinction between Riemannian conceptual-type math vs previous computational, apparently explicit statement from Hilbert, in Zahlbreicht or something)
% Mathematics at Göttingen under the Nazis

\section{[Part 1]}
At the turn of the 20th century, mathematics was undergoing a revolution led by David Hilbert and his colleagues at G\"ottingen. A new style of mathematics was being developed from the ideas of Gauss, Dirichlet, Riemann and Dedekind. Their approach may be roughly characterised with the following traits:
\begin{enumerate}
    \item A preference for abstract definitions.
    \item Complete acceptance of infinite sets.
    \item Conceptual proofs over computational proofs.
    \item Focus on structures defined axiomatically
    \item Reliance on purely existential proofs.
\end{enumerate}
The grand cathedral that was to house their works was the then new Zermelo–Fraenkel set theory. A collection of axioms from which all  mathematics could be derived from. It echoed the ideas of the logicists (i.e. the idea that mathematics can be reduced to ``logical theory''), but designed intentional to ground their mathematicians. The powerful system made strong existential claims about the nature of infinite sets, exemplifying (2) in its wholehearted acceptance of the controversial notions of infinity that Cantor and Dedekind had introduced a decade prior.

\begin{itemize}
    \item More about the "first" foundational crisis, with the paradoxes of set theory, and what not.
    \item If I can, I'd like to cite a chunk of this section away.
\end{itemize}

\subsection{Kronecker}
\begin{itemize}
    \item \emph{"God created the whole numbers, all else is the work of man."}
    \item His relationship with Georg Cantor
    \item His opinion on finitism and problems with Hilbert.
    \item Also a thing on how Hilbert sees him. (foreshadowing, how Hilbert sees Kronecker in Brouwer later).
\end{itemize}

\subsection{Hermann Weyl}

\begin{itemize}
    \item Introduce Weyl as student of Hilbert
    \item Opinions on this new axiomatic, approach and the paradoxes:
    \begin{itemize}
        \item "from a system of intuitive results into a game with formulas that proceeds according to fixed rules"
        \item  "If Hilbert's view prevails over intuitionism, as appears to be the case, then I see in this a decisive defeat of the philosophical attitude of pure phenomenology, which thus proves to be insufficient for the understanding of creative science even in the area of cognition that is most primal and most readily open to evidence – mathematics." 
    \end{itemize}
\end{itemize}

\section{[Part 2] - Brouwer-Hilbert Controversy}

\subsection{L. E. J. Brouwer}

\begin{itemize}
    \item Introduction to Brouwer and his differences with Hilbert
    \item Brouwer being "revolutionary"
    \item Weyl's diagnosis as a "On the New Foundational Crisis in Mathematics"
\end{itemize}

\section{[Part 3] - Hilbert's Program and G\"odel's Incompleteness}

\section{[Part 4] - The aftermath of G\"odel}

\begin{itemize}
    \item Computability theory
    \item BHK-Interpretation
    \item Curry-Howard correspondance
\end{itemize}

\section{[Part 5] - Modern research in Intuitionism}
\begin{itemize}
    \item Connections to programming language semantics
    \item Topological semantics for Intuitionism
\end{itemize}

%(how i feel about the future)
%When we look back at these sentiments of scientism, my first thought is that these people must have been naive. Yet, when I ruminate about the problems the world faces today, I am often overcome by despair. Otherwise, I sometimes find myself hoping that some new magical science or technology will swoop in and solve all these problems that seem so insurmountable. I wonder if those who had witnessed triumphs of the Scientific Revolution of the previous centuries felt similarly.

%https://www.youtube.com/watch?v=ltjI3BXKBgY
%https://www.psychiatrymargins.com/p/emil-du-bois-reymond-and-ignorance
%https://en.wikipedia.org/wiki/Ignoramus_et_ignorabimus
%https://en.wikipedia.org/wiki/Emil_du_Bois-Reymond
%https://ima.org.uk/29760/historical-notes-the-rise-and-fall-of-the-consistency-programme/
%https://www.researchgate.net/publication/286973084_The_crisis_in_the_foundations_of_mathematics
%\begin{chapquote}{Aldous Huxley}
%    Consistency is contrary to nature, contrary to life.
%\end{chapquote}
